\chapter{Tensors, shapes, layers and the arithmetic of convolutions}\label{ch:shapes}

\begin{onething}
A shape such as \shape{(B, N, C)} is a precise statement about how numbers are arranged, and the
three layer types this model is built from (the linear layer, the convolution and the $1\times1$
convolution) each change shapes and own parameters by formulas you can work out by hand before the
computer tells you.
\end{onething}

You are about to read a model whose every line changes the shape of a tensor. Most bugs in tensor code
are shape bugs, and the worst of them do not crash: a wrong reshape produces a tensor of the right size
with its numbers in the wrong places, and training carries on with scrambled data. The defence is to
know, for every operation, what shape goes in, what shape comes out, and in what order the numbers sit.

This chapter gives you that knowledge in four pieces. First the vocabulary: tensors, axes, the batch,
and the two layouts the model keeps switching between. Second the linear layer. Third the
convolution, worked by hand, with the formula for its output size derived by counting and then checked
on every stage of Indradhanu's encoder. Fourth the $1\times1$ convolution, which ties the other two
together and is the subject of \cref{ch:compressor}. Throughout, you will predict parameter counts by
hand and check them against the real model.

\begin{analogy}
Think of a tensor as a warehouse. The numbers are boxes laid out along one very long aisle (computer
memory is one-dimensional). The \emph{shape} is the warehouse's filing system: ``aisle positions
0--15 are channel~0, 16--31 are channel~1'', and within a channel ``four rows of four''. Two
warehouses can hold exactly the same boxes in exactly the same aisle order and still mean different
things if their filing systems differ. A \emph{reshape} changes the filing system without moving a box.
A \emph{transpose} changes which boxes the filing system points at. If you relabel the aisle when you
needed to move the boxes, every box is filed under the wrong name, and nothing tells you.
\end{analogy}

%======================================================================
\section{Tensors and how to read a shape}\label{sec:shapes:tensors}

\begin{definition}[Tensor, axis, shape]
A \emph{tensor}\index{tensor} is a rectangular block of numbers. Each direction along which you can
index it is an \emph{axis}\index{axis} (also called a dimension), and the list of the lengths of its
axes is its \emph{shape}\index{shape}. The number of axes is the tensor's \emph{rank}\index{rank (of a
tensor)}.
\end{definition}

\begin{center}
\begin{tabular}{lll}
\toprule
rank & example shape & what it is\\
\midrule
0 & \shape{()} & a single number\\
1 & \shape{(5,)} & a list of five numbers\\
2 & \shape{(3, 4)} & a grid, three rows by four columns\\
3 & \shape{(165, 128, 128)} & a stack of 165 grids: one hyperspectral patch\\
4 & \shape{(8, 165, 128, 128)} & eight such patches\\
\bottomrule
\end{tabular}
\end{center}

Read a shape left to right as nested containers. \shape{(8, 165, 128, 128)} is ``eight items, each
containing 165 layers, each of which is 128 rows by 128 columns''. The number of values in a tensor is
the product of its shape: $8\times165\times128\times128 = 21{,}626{,}880$.

\subsection{The batch}

Neural networks process several examples at once, because a graphics processor is far more efficient
on 128 patches together than on one patch 128 times. That group is a \emph{batch}\index{batch}, and its
size is written $B$. Almost every tensor in the model has $B$ as its first axis, and almost every
operation treats the $B$ items independently. When you read the code you can usually ignore $B$ and
think about one patch.

\subsection{The symbols used in this book}

\begin{center}
\begin{tabular}{lll}
\toprule
symbol & meaning & typical value\\
\midrule
$B$ & batch size & 128 in training$^{\ast}$\\
$C$ & input spectral bands & 165\\
$D$ & channels after the spectral compressor & 32 (24 in v0.1.0)\\
$H$, $W$ & height and width in pixels & 128, 128\\
$N$ & number of tokens at a stage & 1024 at stage 1\\
$C_i$ & token width at encoder stage $i$ & 32, 64, 160, 256\\
$k$, $s$, $p$ & kernel, stride, padding of a convolution & \cref{sec:shapes:conv}\\
\bottomrule
\end{tabular}
\end{center}
{\small $^{\ast}$Batch size 128 is the v0.2.0 training
configuration\tagM{configs/hyperspectral_segformer_exp_2.json}; the code default is 32.}

Watch the two meanings of ``channels''. $C=165$ is the number of input bands. $C_i$ is the width of
the vectors inside encoder stage $i$. They are unrelated numbers that both get called channels.

\subsection{Where the numbers live: row-major order}\label{sec:shapes:memory}

Computer memory is a single line of numbers. A tensor of shape \shape{(B, C, H, W)} is stored in
\emph{row-major order}\index{row-major order}: the last axis changes fastest. Walking along memory, you
pass $W$ values of row~0 of channel~0, then row~1 of channel~0, and so on through all $H$ rows; then
the same for channel~1; and after all $C$ channels, the next batch item.

You can derive the memory position (the \emph{offset}) of element $(b,c,h,w)$, counting everything
from $0$, by counting how many values come before it:
\begin{align*}
\text{offset}(b,c,h,w)
&= b\cdot(CHW) + c\cdot(HW) + h\cdot W + w
  && \text{items, channels, rows, columns before it}\\
&= \bigl((b\,C + c)\,H + h\bigr)\,W + w.
  && \text{factor out $W$, then $H$}
\end{align*}
The numbers $CHW$, $HW$, $W$ and $1$ multiplying the indices are called the tensor's
\emph{strides}\index{stride!of a tensor} (not to be confused with a convolution's stride in
\cref{sec:shapes:conv}). They say how far apart in memory two neighbours along each axis are.

\begin{worked}[label={we:shapes:offset}]{Finding a value in memory}
Take a tensor of shape \shape{(1, 2, 2, 2)}: one item, two channels, each a $2\times2$ grid. Where is
element $(b,c,h,w) = (0,1,0,1)$, that is channel~1, row~0, column~1? The strides are
$CHW = 2\cdot2\cdot2 = 8$, $HW = 2\cdot2 = 4$, $W = 2$ and $1$. So
\begin{align*}
\text{offset} &= 0\cdot8 + 1\cdot4 + 0\cdot2 + 1\cdot1\\
&= 0 + 4 + 0 + 1\\
&= 5.
\end{align*}
It is the sixth value in memory (offset $5$, counting from $0$): four values of channel~0 come first,
then channel~1's row~0, column~0 at offset~4, then our element.
\end{worked}

%======================================================================
\section{Two layouts of the same numbers}\label{sec:shapes:layouts}

The model keeps its data in one of two arrangements.

\begin{center}
\begin{tabular}{lll}
\toprule
layout & shape & used by\\
\midrule
spatial\index{layout!spatial} & \shape{(B, C, H, W)} & convolutions, pixel shuffle, the loss\\
sequence\index{layout!sequence} & \shape{(B, N, C)} & attention, LayerNorm, linear layers\\
\bottomrule
\end{tabular}
\end{center}

Spatial says ``a picture with $C$ layers''. Sequence says ``a list of $N$ \emph{tokens}\index{token},
each described by $C$ numbers''. A token here is one spatial position's vector of channel values; when
you convert spatial to sequence, $N = H\times W$. The conversion is written in the code as
\begin{lstlisting}
x = x.flatten(2)        # (B, C, H, W) -> (B, C, H*W) = (B, C, N)
x = x.transpose(1, 2)   # (B, C, N)    -> (B, N, C)
\end{lstlisting}
\code{flatten(2)} merges every axis from axis~2 onwards into one: the two spatial axes become a single
list of $N$ positions. \code{transpose(1, 2)} swaps axes~1 and~2. The way back is
\begin{lstlisting}
x = x.transpose(1, 2).reshape(B, C, H, W)   # (B, N, C) -> (B, C, N) -> (B, C, H, W)
\end{lstlisting}

\coursefig[\linewidth]{iw-07-layouts}{The two layouts of the same numbers. (a) Spatial
\texttt{(B, C, H, W)}: a grid of positions per channel, as convolutions need. (b) \texttt{flatten(2)}
joins the two spatial axes into $N = H\times W$ positions. (c) \texttt{transpose(1, 2)} makes each
position a row of $C$ numbers, a token, as attention and LayerNorm need. A reshape alone would scramble
the order. (Course figure IW~7.1.)}{fig:shapes:layouts}

\subsection{Why a reshape alone scrambles the data}

A \code{reshape} changes the shape but keeps the numbers in the same memory order. The two layouts need
\emph{different} memory orders: spatial wants all of channel~0's pixels, then all of channel~1's;
sequence wants token~0's channel values, then token~1's. So an axis swap is unavoidable. The smallest
example that shows this has two channels and four pixels.

\begin{worked}[label={we:shapes:layout}]{A $2\times2\times2$ cube through flatten, transpose and a wrong reshape}
Take $B=1$, $C=2$, $H=W=2$, holding the numbers $1$ to $8$ in memory order:
\[
\text{channel 0} = \begin{pmatrix}1&2\\3&4\end{pmatrix},\qquad
\text{channel 1} = \begin{pmatrix}5&6\\7&8\end{pmatrix},\qquad
\text{memory: } 1\ 2\ 3\ 4\ 5\ 6\ 7\ 8 .
\]
The four pixels are: top-left (channel values $1$ and $5$), top-right ($2$, $6$), bottom-left ($3$,
$7$) and bottom-right ($4$, $8$). The correct tokens are therefore $[1,5]$, $[2,6]$, $[3,7]$, $[4,8]$.

\emph{Step 1, flatten(2).} Shape \shape{(1, 2, 2, 2)} becomes \shape{(1, 2, 4)}. Each channel's
$2\times2$ grid is read row by row into a list of four:
\[
\begin{pmatrix}1&2&3&4\\5&6&7&8\end{pmatrix}.
\]
Memory has not moved: still $1\,2\,3\,4\,5\,6\,7\,8$.

\emph{Step 2, transpose(1, 2).} Shape \shape{(1, 2, 4)} becomes \shape{(1, 4, 2)}: rows and columns
swap.
\[
\begin{pmatrix}1&5\\2&6\\3&7\\4&8\end{pmatrix}.
\]
Row~$n$ is token~$n$, and each token holds its pixel's two channel values. Correct. Laid out
contiguously, memory now reads $1\,5\,2\,6\,3\,7\,4\,8$: the boxes have moved.

\emph{The wrong way, reshape(1, 4, 2) straight from the cube.} A reshape reads memory in order,
$1\,2\,3\,4\,5\,6\,7\,8$, and cuts it into four rows of two:
\[
\begin{pmatrix}1&2\\3&4\\5&6\\7&8\end{pmatrix}.
\]
The ``token'' $[1,2]$ is two different pixels' channel-0 values. Every token is wrong, the shape is
right, and no error is raised.
\end{worked}

The source carries the same warning, in \file{app/foundation_models/components/overlap_patch_embedding.py}:
``Cannot use view/reshape directly because the axes need to swap''.

\subsection{Broadcasting}

When two tensors of different shapes are multiplied, numpy and torch stretch any axis of length~1 to
match the other tensor. This is \emph{broadcasting}\index{broadcasting}:
\[
\underbrace{\shape{(B, 165, H, W)}}_{\text{the cube}} \times
\underbrace{\shape{(B, 1, H, W)}}_{\text{a pixel mask}} \to \shape{(B, 165, H, W)}.
\]
The mask's single channel is reused for all 165 bands. That is why pixel masks in this code have shape
\shape{(B, 1, H, W)} and not \shape{(B, H, W)}: the $1$ is a placeholder in the channel slot, so the
axes line up. In the same way the per-band normalisation statistics are stored with shape
\shape{(1, 165, 1, 1)} (\file{app/foundation_models/components/pixel_normalization.py}) so that one
mean per band is reused at every pixel of every patch.

%======================================================================
\section{The linear layer}\label{sec:shapes:linear}

\subsection{One neuron}

A \emph{neuron}\index{neuron} takes a vector $\vect{x}\in\R^{n}$, multiplies each entry by its own
\emph{weight}\index{weight}, adds the products, and adds a \emph{bias}\index{bias}:
\begin{equation}\label{eq:shapes:neuron}
y = w_1x_1 + w_2x_2 + \dots + w_nx_n + b = \vect{w}\T\vect{x} + b .
\end{equation}
The sum of products $\vect{w}\T\vect{x}$ is the dot product (\cref{ch:linalg}). A neuron has $n$
weights and one bias: $n+1$ parameters. In most networks a nonlinear function (an
\emph{activation}\index{activation function}, such as GELU in \cref{ch:block}) is applied to $y$
afterwards; here we count only the linear part, because that is where the parameters are.

\subsection{A layer of neurons is a matrix}

Put $m$ neurons side by side, all reading the same $\vect{x}$. Stack their weight vectors as the rows
of a matrix $\mat{W}\in\R^{m\times n}$ and their biases into $\vect{b}\in\R^{m}$:
\begin{equation}\label{eq:shapes:linear}
\vect{y} = \mat{W}\vect{x} + \vect{b},\qquad
y_i = \sum_{j=1}^{n} W_{ij}\,x_j + b_i \quad (i = 1,\dots,m).
\end{equation}
This is a \emph{linear layer}\index{linear layer} (PyTorch \code{nn.Linear(n, m)}). Output $i$ is a
weighted mixture of all the inputs, with its own recipe of weights, row $i$ of $\mat{W}$. Counting:
\begin{equation}\label{eq:shapes:linearcount}
\text{parameters} = \underbrace{m\cdot n}_{\text{entries of } \mat{W}} + \underbrace{m}_{\text{entries of } \vect{b}}
= \text{in}\times\text{out} + \text{out}.
\end{equation}

\begin{worked}[label={we:shapes:linear}]{A linear layer from three inputs to two outputs}
Let $n=3$, $m=2$,
\[
\mat{W} = \begin{pmatrix}1.0 & 0.0 & 0.0\\ 0.0 & 0.5 & 0.5\end{pmatrix},\quad
\vect{b} = \begin{pmatrix}0.0\\0.1\end{pmatrix},\quad
\vect{x} = \begin{pmatrix}0.10\\0.20\\0.30\end{pmatrix}.
\]
Row~1 of $\mat{W}$ is $(1.0, 0.0, 0.0)$:
\begin{align*}
y_1 &= 1.0\times0.10 + 0.0\times0.20 + 0.0\times0.30 + 0.0\\
&= 0.10 + 0 + 0 + 0 = 0.10.
\end{align*}
Row~2 is $(0.0, 0.5, 0.5)$:
\begin{align*}
y_2 &= 0.0\times0.10 + 0.5\times0.20 + 0.5\times0.30 + 0.1\\
&= 0 + 0.10 + 0.15 + 0.1 = 0.35.
\end{align*}
So $\vect{y} = (0.10,\ 0.35)\T$. Output~1 copied input~1; output~2 averaged inputs~2 and~3 and added
an offset. Parameter count: $3\times2 + 2 = 8$, which is what PyTorch reports for
\code{nn.Linear(3, 2)}\tagC.
\end{worked}

\subsection{Applied to a sequence}

\code{nn.Linear} acts on the \emph{last} axis of its input and treats all other axes as a batch. Given
tokens of shape \shape{(B, N, C\_in)} it applies the same $\mat{W}$ and $\vect{b}$ to every one of the
$B\cdot N$ tokens and returns \shape{(B, N, C\_out)}. That is why linear layers live in the sequence
layout: the last axis must be the feature axis.

\subsection{Why not one linear layer for the whole image?}

You could flatten an entire patch into one vector and apply one linear layer to it. Count the cost of
mapping a patch to a patch of the same size: the input has $165\times128\times128 = 2{,}703{,}360$
values, so
\[
\text{weights} = 2{,}703{,}360 \times 2{,}703{,}360 \approx 7.3\times10^{12}\tagC .
\]
Seven trillion parameters, for one layer. Worse, such a layer treats a field in the top-left corner and
the same field in the bottom-right as unrelated problems. The convolution fixes both faults.

%======================================================================
\section{The convolution}\label{sec:shapes:conv}

\subsection{Two ideas: locality and sharing}

A convolution builds each output value from a small neighbourhood of the input
(\emph{locality}\index{locality}), and it uses the same weights at every position
(\emph{weight sharing}\index{weight sharing}). The small block of weights is the \emph{kernel}%
\index{kernel}; its side length is the \emph{kernel size} $k$. The kernel slides over the input; at
each position it multiplies the values under it by its weights and adds them up.

\subsection{Single channel, written as a sum}

\begin{definition}[2-D cross-correlation]
For a single-channel input $x$ (an $H\times W$ grid) and a $k\times k$ kernel $K$ with bias $b$, the
output at row $i$, column $j$ (counting from $0$) is
\begin{equation}\label{eq:shapes:xcorr}
y_{i,j} = b + \sum_{u=0}^{k-1}\sum_{v=0}^{k-1} K_{u,v}\; x_{i+u,\ j+v}.
\end{equation}
\end{definition}
Deep-learning libraries call this a \emph{convolution}\index{convolution}. Strictly it is a
\emph{cross-correlation}\index{cross-correlation}: the mathematical convolution flips the kernel first
($K_{k-1-u,\,k-1-v}$). Because the kernel is learned, the flip makes no difference to what a network can
represent, and this book follows the library usage.

Read \cref{eq:shapes:xcorr} as ``place the kernel's top-left corner on input position $(i,j)$,
multiply the overlapping pairs, add''. With an $H\times W$ input and no padding the corner can go at
rows $0$ to $H-k$ and columns $0$ to $W-k$, so the output is $(H-k+1)\times(W-k+1)$.

\begin{worked}[label={we:shapes:conv}]{A $3\times3$ input and a $2\times2$ kernel, by hand}
Do this on paper. Input and kernel, bias $b=0$:
\[
x = \begin{pmatrix}1&2&0\\0&1&3\\2&1&1\end{pmatrix},\qquad
K = \begin{pmatrix}1&0\\-1&2\end{pmatrix}.
\]
The output is $(3-2+1)\times(3-2+1) = 2\times2$. Each entry pairs $K$ with the $2\times2$ block of $x$
whose top-left corner is at $(i,j)$.

$(i,j)=(0,0)$, block $\begin{pmatrix}1&2\\0&1\end{pmatrix}$:
\begin{align*}
y_{0,0} &= 1\times1 + 0\times2 + (-1)\times0 + 2\times1\\
&= 1 + 0 + 0 + 2 = 3.
\end{align*}
$(0,1)$, block $\begin{pmatrix}2&0\\1&3\end{pmatrix}$:
\begin{align*}
y_{0,1} &= 1\times2 + 0\times0 + (-1)\times1 + 2\times3\\
&= 2 + 0 - 1 + 6 = 7.
\end{align*}
$(1,0)$, block $\begin{pmatrix}0&1\\2&1\end{pmatrix}$:
\begin{align*}
y_{1,0} &= 1\times0 + 0\times1 + (-1)\times2 + 2\times1\\
&= 0 + 0 - 2 + 2 = 0.
\end{align*}
$(1,1)$, block $\begin{pmatrix}1&3\\1&1\end{pmatrix}$:
\begin{align*}
y_{1,1} &= 1\times1 + 0\times3 + (-1)\times1 + 2\times1\\
&= 1 + 0 - 1 + 2 = 2.
\end{align*}
So
\[
y = \begin{pmatrix}3&7\\0&2\end{pmatrix}.
\]
PyTorch's \code{F.conv2d} gives the same four numbers\tagC. This kernel responds most to ``top-left and
bottom-right bright, bottom-left dark'', and its four weights are used at every position.
\end{worked}

\subsection{Stride and padding}

Two more settings control where the kernel goes.

\begin{definition}[Stride and padding]
The \emph{stride}\index{stride!of a convolution} $s$ is how far the kernel jumps between positions:
with $s=2$ it is placed at every second row and every second column. The \emph{padding}\index{padding}
$p$ is a border of $p$ zeros added on every side of the input before the kernel slides.
\end{definition}

With stride and padding, \cref{eq:shapes:xcorr} becomes
\begin{equation}\label{eq:shapes:xcorrsp}
y_{i,j} = b + \sum_{u=0}^{k-1}\sum_{v=0}^{k-1} K_{u,v}\; \tilde{x}_{s i+u,\ s j+v},
\end{equation}
where $\tilde{x}$ is the zero-padded input (so $\tilde{x}_{a,b'} = x_{a-p,\,b'-p}$ inside the original
picture and $0$ in the border). The kernel's top-left corner sits at padded position $(si, sj)$.

\begin{figure}[htbp]\centering
\begin{tikzpicture}[scale=0.62]
% padded 6x6 grid: border cells are padding
\foreach \r in {0,...,5}{\foreach \c in {0,...,5}{
  \pgfmathsetmacro{\pad}{(\r==0 || \r==5 || \c==0 || \c==5) ? 1 : 0}
  \ifdim\pad pt=1pt
    \draw[maskgrey,dashed,fill=maskgrey!12] (\c,-\r) rectangle ++(1,-1);
    \node[font=\scriptsize,maskgrey] at (\c+0.5,-\r-0.5) {0};
  \else
    \draw[spat] (\c,-\r) rectangle ++(1,-1);
  \fi}}
\foreach \r in {1,...,4}{\foreach \c in {1,...,4}{
  \pgfmathtruncatemacro{\idx}{(\r-1)*4+\c-1}
  \node[font=\scriptsize] at (\c+0.5,-\r-0.5) {$x_{\idx}$};}}
% window at output (0,0): padded rows 0-2, cols 0-2
\draw[loss,very thick] (0,0) rectangle (3,-3);
% window at output (0,1): padded cols 2-4
\draw[attn,very thick,dashed] (2,0.08) rectangle (5,-3.08);
\node[font=\small,anchor=west] at (6.4,-0.6) {\textcolor{loss}{solid}: window for $y_{0,0}$, corner at padded $(0,0)$};
\node[font=\small,anchor=west] at (6.4,-1.5) {\textcolor{attn}{dashed}: window for $y_{0,1}$, corner at padded $(0,2)$};
\node[font=\small,anchor=west] at (6.4,-2.4) {grey border: padding $p=1$ of zeros};
\node[font=\small,anchor=west] at (6.4,-3.3) {kernel $k=3$, stride $s=2$};
\node[font=\small,anchor=west] at (6.4,-4.4) {output $2\times2$; real pixels seen:};
\node[font=\small,anchor=west] at (6.4,-5.2) {$y_{0,0}$: 4, $y_{0,1}$: 6, $y_{1,0}$: 6, $y_{1,1}$: 9};
\end{tikzpicture}
\caption{A $3\times3$ kernel sliding with stride~2 over a $4\times4$ input padded by one ring of zeros,
the setting of encoder stages~2 to~4. Consecutive windows share a column (the overlap is $k-s=1$). The
corner windows see only 4 real pixels; the bottom-right one sees 9.}
\label{fig:shapes:slide}
\end{figure}

\subsection{Many channels in, many channels out}

Real inputs have $C_{\text{in}}$ channels, and a layer produces $C_{\text{out}}$ output channels. Each
output channel $o$ has its own kernel for every input channel, so the weights form a four-axis tensor
$K\in\R^{C_{\text{out}}\times C_{\text{in}}\times k\times k}$, and each output channel has one bias:
\begin{equation}\label{eq:shapes:multiconv}
y_{o,i,j} = b_o + \sum_{c=0}^{C_{\text{in}}-1}\ \sum_{u=0}^{k-1}\sum_{v=0}^{k-1}
K_{o,c,u,v}\; \tilde{x}_{c,\ s i+u,\ s j+v}.
\end{equation}
In words: output channel $o$ at position $(i,j)$ looks at a $k\times k$ window in \emph{every} input
channel, and mixes all $C_{\text{in}}\cdot k\cdot k$ values with its own weights. The parameter count is
\begin{equation}\label{eq:shapes:convcount}
\text{parameters} = C_{\text{in}}\cdot C_{\text{out}}\cdot k\cdot k + C_{\text{out}}.
\end{equation}
For example \code{nn.Conv2d(3, 2, kernel_size=2)} has $3\cdot2\cdot2\cdot2 + 2 = 24 + 2 = 26$
parameters\tagC. Notice what is \emph{absent} from \cref{eq:shapes:convcount}: $H$ and $W$. Weight
sharing makes the cost independent of the picture size.

%======================================================================
\section{The output-size formula, derived by counting}\label{sec:shapes:outsize}

You will use this formula more than any other in Part~III, so derive it rather than memorise it.

Work along one axis; the other is the same. After padding, the axis has length
\[
L = H + 2p .
\]
The kernel's first cell sits at positions $0, s, 2s, 3s, \dots$, that is at $ns$ for $n = 0, 1, 2,
\dots$. Placed at $ns$, the kernel covers positions $ns$ to $ns + k - 1$. It must fit inside the padded
axis, so its last cell must be at most $L-1$:
\begin{align*}
ns + k - 1 &\le L - 1
  && \text{the last covered cell is inside}\\
ns &\le L - k
  && \text{add $1-k$ to both sides}\\
n &\le \frac{L-k}{s}.
  && \text{divide by $s>0$}
\end{align*}
Since $n$ is a whole number, the largest allowed $n$ is $\left\lfloor (L-k)/s \right\rfloor$, where
$\lfloor\cdot\rfloor$ (the \emph{floor}\index{floor}) means ``round down to a whole number''. The
allowed values are $n = 0, 1, \dots, \lfloor (L-k)/s\rfloor$; counting the $0$ as well, there are
\begin{equation}\label{eq:shapes:outsize}
H_{\text{out}} = \left\lfloor \frac{H + 2p - k}{s} \right\rfloor + 1
\end{equation}
of them. Each kernel position gives one output value, so this is the output size.

\begin{worked}[label={we:shapes:stages}]{Every stage of Indradhanu's encoder}
The encoder builds its four patch embeddings in
\file{app/foundation_models/components/seg_former_encoder.py} (lines~110--114):
stage~1 has $k=4$, $s=4$, $p=0$; stages~2 to~4 have $k=3$, $s=2$, $p=1$. Input $H=128$.

Stage 1:
\begin{align*}
H + 2p - k &= 128 + 0 - 4 = 124,\\
124 / 4 &= 31, \quad \lfloor 31 \rfloor = 31,\\
H_{\text{out}} &= 31 + 1 = 32.
\end{align*}
Stage 2 ($H=32$):
\begin{align*}
32 + 2 - 3 &= 31,\\
31 / 2 &= 15.5, \quad \lfloor 15.5\rfloor = 15,\\
H_{\text{out}} &= 15 + 1 = 16.
\end{align*}
Stage 3 ($H=16$):
\begin{align*}
16 + 2 - 3 &= 15,\\
15/2 &= 7.5, \quad \lfloor 7.5 \rfloor = 7,\\
H_{\text{out}} &= 7 + 1 = 8.
\end{align*}
Stage 4 ($H=8$):
\begin{align*}
8 + 2 - 3 &= 7,\\
7/2 &= 3.5,\quad \lfloor 3.5\rfloor = 3,\\
H_{\text{out}} &= 3 + 1 = 4.
\end{align*}
The grids are $32, 16, 8, 4$ on a side, so the token counts are $32^2 = 1024$, $16^2 = 256$,
$8^2 = 64$ and $4^2 = 16$. Running the real encoder on a \shape{(1, 32, 128, 128)} input returns
feature maps \shape{(1, 32, 32, 32)}, \shape{(1, 64, 16, 16)}, \shape{(1, 160, 8, 8)} and
\shape{(1, 256, 4, 4)}\tagC.
\end{worked}

Two small facts drop out of \cref{eq:shapes:outsize}. When $k = s$ and $p = 0$, it gives
$\lfloor (H-k)/k\rfloor + 1 = \lfloor H/k \rfloor$: the windows tile the axis. When $k=3$, $s=2$,
$p=1$, it gives $\lfloor (H-1)/2 \rfloor + 1$, which equals $H/2$ when $H$ is even.

\begin{worked}[label={we:shapes:pad}]{Why stage 1 must set its padding explicitly}
\code{OverlapPatchEmbedding} defaults to padding $p = \lfloor k/2 \rfloor$ when none is given. Had stage~1
used that default with $k=4$, the padding would be $2$:
\begin{align*}
128 + 2\times2 - 4 &= 128 + 4 - 4 = 128,\\
128 / 4 &= 32,\\
H_{\text{out}} &= 32 + 1 = 33.
\end{align*}
A $33\times33$ grid of $1089$ tokens, windows hanging two pixels over the edge, and token boundaries
no longer aligned with the $4\times4$ blocks that the masking code assumes. The encoder therefore
passes \lstinline!padding = 0 if i == 0 else None!.
\end{worked}

\begin{sourcenote}[: ``H\_out = H\_in // stride'' holds only for even sizes]
The docstring of \code{OverlapPatchEmbedding}
(\file{app/foundation_models/components/overlap_patch_embedding.py}, lines~87--105) states
``\code{H' = H // desired_compression}'' and ``with padding = patch\_size // 2:
\code{H_out = H_in // stride}''. \Cref{eq:shapes:outsize} says otherwise for odd inputs. With $k=3$,
$s=2$, $p=1$ and $H=33$: $\lfloor(33+2-3)/2\rfloor+1 = \lfloor 16\rfloor + 1 = 17$, while
$33 /\!/ 2 = 16$. PyTorch returns 17\tagC. The docstring line~87 also says the convolution uses
\code{padding=patch_size//2}, which is false for stage~1 ($p=0$). Neither matters for the $128\times128$
patches the model is trained on, because every stage then receives an even size; the formula is the
authority.
\end{sourcenote}

\begin{worked}[label={we:shapes:spatialenc}]{The convolutional autoencoder and a $100\times100$ patch}
The old \code{SpatialEncoder} (\cref{ch:recon}) uses $k=4$, $s=2$, $p=1$ at each of three stages. On
$H=128$:
\[
\lfloor(128 + 2 - 4)/2\rfloor + 1 = \lfloor 126/2 \rfloor + 1 = 63 + 1 = 64,
\]
an exact halving; then $32$, then $16$. On $H = 100$:
\begin{align*}
\text{stage 1: } & \lfloor (100 + 2 - 4)/2 \rfloor + 1 = \lfloor 98/2 \rfloor + 1 = 49 + 1 = 50,\\
\text{stage 2: } & \lfloor (50 + 2 - 4)/2 \rfloor + 1 = \lfloor 48/2 \rfloor + 1 = 24 + 1 = 25,\\
\text{stage 3: } & \lfloor (25 + 2 - 4)/2 \rfloor + 1 = \lfloor 23/2\rfloor + 1 = 11 + 1 = 12.
\end{align*}
The decoder doubles three times, $12\to24\to48\to96$, and $96\ne100$: the output cannot be subtracted
from the input. Height and width must be divisible by $2^3 = 8$.
\end{worked}

%======================================================================
\section{Receptive field}\label{sec:shapes:rf}

\begin{definition}[Receptive field]
The \emph{receptive field}\index{receptive field} of an output value is the set of input pixels that can
affect it. For a stack of convolutions it is a square; its side is what this book calls the receptive
field size.
\end{definition}

One convolution with kernel $k$ has a receptive field of $k$ pixels. Stack layers and it grows. Keep
two numbers per layer $\ell$: the receptive field $r_\ell$ and the \emph{jump}\index{jump} $J_\ell$, the
distance in input pixels between two neighbouring outputs of layer $\ell$. Start from the input itself:
$r_0 = 1$, $J_0 = 1$.

Layer $\ell$ has kernel $k_\ell$ and stride $s_\ell$. Its window spans $k_\ell$ neighbouring outputs of
layer $\ell-1$. Those are $J_{\ell-1}$ input pixels apart, so the window's first and last positions are
$(k_\ell-1)J_{\ell-1}$ input pixels apart, and each of them brings its own receptive field of
$r_{\ell-1}$:
\begin{align}
r_\ell &= r_{\ell-1} + (k_\ell - 1)\,J_{\ell-1}, \label{eq:shapes:rf}\\
J_\ell &= J_{\ell-1}\, s_\ell. \label{eq:shapes:jump}
\end{align}
The jump multiplies because moving one output of layer $\ell$ moves $s_\ell$ outputs of layer
$\ell-1$, each $J_{\ell-1}$ pixels.

\begin{worked}[label={we:shapes:rf}]{Receptive field of the four patch embeddings}
Stage~1, $k=4$, $s=4$:
\begin{align*}
r_1 &= 1 + (4-1)\times1 = 1 + 3 = 4, & J_1 &= 1\times4 = 4.
\end{align*}
Stage~2, $k=3$, $s=2$:
\begin{align*}
r_2 &= 4 + (3-1)\times4 = 4 + 8 = 12, & J_2 &= 4\times2 = 8.
\end{align*}
Stage~3:
\begin{align*}
r_3 &= 12 + 2\times8 = 12 + 16 = 28, & J_3 &= 8\times2 = 16.
\end{align*}
Stage~4:
\begin{align*}
r_4 &= 28 + 2\times16 = 28 + 32 = 60, & J_4 &= 16\times2 = 32.
\end{align*}
A stage-4 token sits on a $32$-pixel grid ($J_4$) but its patch-embedding path sees a $60\times60$
window of the input ($r_4$). Back-propagating from one centre output through four such convolutions
with all-one weights touches exactly 4, 12, 28 and 60 input rows\tagC.
\end{worked}

Two cautions. First, these numbers count only the patch-embedding convolutions. The transformer blocks
in between contain a $3\times3$ depthwise convolution (\cref{ch:block}) that widens the field further,
and attention (\cref{ch:attention}) lets a token read every other token at its stage, which makes the
effective receptive field the whole patch. Second, ``a stage-4 token covers a $32\times32$ region'' (as
the Indradhanu course puts it) describes the jump: the grid cell each token stands for. Its receptive
field is wider, because the stride-2 windows of stages~2 to~4 overlap.

%======================================================================
\section{The $1\times1$ convolution is a linear layer at every pixel}\label{sec:shapes:onebyone}

Set $k=1$, $s=1$, $p=0$ in \cref{eq:shapes:multiconv}. The sums over $u$ and $v$ have a single term
($u=v=0$), the padded input is the input, and the kernel tensor $K_{o,c,0,0}$ is a matrix $W_{o,c}$:
\begin{align}
y_{o,i,j} &= b_o + \sum_{c} K_{o,c,0,0}\; x_{c,\,i,\,j}
  && \text{only $u=v=0$ survives}\notag\\
&= b_o + \sum_{c} W_{o,c}\; x_{c,i,j}.
  && \text{rename $K_{o,c,0,0}$ as $W_{o,c}$}\label{eq:shapes:onebyone}
\end{align}
Compare with \cref{eq:shapes:linear}. At each pixel $(i,j)$, the vector of channel values
$\vect{x}_{:,i,j}$ is multiplied by the same matrix $\mat{W}$ and shifted by the same $\vect{b}$:
\begin{equation}
\vect{y}_{:,i,j} = \mat{W}\,\vect{x}_{:,i,j} + \vect{b}\qquad\text{for every pixel } (i,j).
\end{equation}
A \emph{$1\times1$ convolution}\index{convolution!$1\times1$} is a linear layer applied independently at
every pixel. It mixes channels and does no spatial mixing at all: pixel $(3,7)$ never looks at
pixel $(3,8)$. Its parameter count, from \cref{eq:shapes:convcount} with $k=1$, is
$C_{\text{in}}\cdot C_{\text{out}} + C_{\text{out}}$, the same as \cref{eq:shapes:linearcount}.

The equivalence is exact in code too: a \code{nn.Conv2d(3, 2, 1)} and a \code{nn.Linear(3, 2)} with
the same weights give identical results, one in the spatial layout and one on the flattened tokens
(checked numerically\tagC).

\coursefig[0.92\linewidth]{iw-08-conv1x1}{A $1\times1$ convolution is a matrix-vector product per pixel.
At every pixel the 4-band vector $\vect{x}$ is multiplied by the same $2\times4$ matrix $\mat{W}$,
giving two output channels, $y_0 = 1.1$ and $y_1 = -0.5$ (the numbers of \cref{we:compressor:mix}).
The spatial grid is untouched; only the channel axis changes. (Course figure IW~8.1.)}{fig:shapes:conv1x1}

\Cref{ch:compressor} uses exactly this: Indradhanu's spectral compressor is a $1\times1$ convolution
from 165 channels to 32.

%======================================================================
\section{Counting parameters: a convolutional encoder by hand}\label{sec:shapes:count}

The best test of understanding is to predict a real model's parameter count before asking the
computer. Take the convolutional autoencoder's encoder of \cref{ch:recon},
\code{SpatialEncoder(in_channels=1, base_channels=32, num_stages=3)}. Each stage is a
\code{Conv2d} with $k=4$, $s=2$, $p=1$, followed by a \code{BatchNorm2d}, which has $2$ parameters per
channel (\cref{ch:compressor} explains why). Channels go $1\to32\to64\to128$.

\begin{worked}[label={we:shapes:spatialcount}]{The SpatialEncoder's 165,024 parameters}
Stage~0, $1\to32$:
\begin{align*}
\text{conv weights} &= 1\times32\times4\times4 = 32\times16 = 512,\\
\text{conv biases} &= 32,\\
\text{BatchNorm} &= 2\times32 = 64,\\
\text{stage total} &= 512 + 32 + 64 = 608.
\end{align*}
Stage~1, $32\to64$:
\begin{align*}
\text{conv weights} &= 32\times64\times16 = 2048\times16 = 32{,}768,\\
\text{conv biases} &= 64,\\
\text{BatchNorm} &= 2\times64 = 128,\\
\text{stage total} &= 32{,}768 + 64 + 128 = 32{,}960.
\end{align*}
Stage~2, $64\to128$:
\begin{align*}
\text{conv weights} &= 64\times128\times16 = 8192\times16 = 131{,}072,\\
\text{conv biases} &= 128,\\
\text{BatchNorm} &= 2\times128 = 256,\\
\text{stage total} &= 131{,}072 + 128 + 256 = 131{,}456.
\end{align*}
Total:
\begin{align*}
608 + 32{,}960 &= 33{,}568,\\
33{,}568 + 131{,}456 &= 165{,}024.
\end{align*}
\lstinline!sum(p.numel() for p in SpatialEncoder(1, 32, 3).parameters())! returns
165{,}024\tagC. The whole autoencoder with its mirrored decoder has
329{,}665\tagM{docs/training/anomaly-detection/10-reconstruction-models.md}.
\end{worked}

%======================================================================
\section{The whole forward pass as a table of shapes}\label{sec:shapes:table}

\Cref{tab:shapes:forward} lists every shape a training patch passes through in Indradhanu, with
$C=165$, $D=32$ and $H=W=128$. You are not expected to understand every row yet. You are expected to
know the table exists, to see that the model starts and ends at \shape{(B, 165, 128, 128)}, and to
check each patch-embedding row against \cref{we:shapes:stages}.

\begin{table}[htbp]\centering\small
\caption{Shapes through one Indradhanu forward pass (course part IW07, with stage rows checked
against the real encoder by \texttt{tools/compute\_shapes.py}). ``$\sim$358'' is the number of
visible stage-1 tokens at mask ratio 0.65 ($1024\times0.35 \approx 358$).}
\label{tab:shapes:forward}
\begin{tabular}{llll}
\toprule
step & module & output shape & chapter\\
\midrule
input & & \shape{(B, 165, 128, 128)} & \\
normalise & \texttt{PixelNormalize} & \shape{(B, 165, 128, 128)} & \ref{ch:bandpipeline}\\
compress & \texttt{SpectralCompressor} & \shape{(B, 32, 128, 128)} & \ref{ch:compressor}\\
S1 patch embed & $k4,s4,p0$ & \shape{(B, 1024, 32)} & \ref{ch:patchembed}\\
S1 delete hidden & \texttt{remove\_tokens} & \shape{(B, $\sim$358, 32)} & \ref{ch:masking}\\
S1 blocks $\times2$ & \texttt{SegFormerBlock} & \shape{(B, $\sim$358, 32)} & \ref{ch:block}\\
S1 restore, to 2-D & & \shape{(B, 32, 32, 32)} = F1 & \ref{ch:encoder}\\
S2 embed + blocks & $k3,s2,p1$ & \shape{(B, 256, 64)} $\to$ F2 \shape{(B, 64, 16, 16)} & \ref{ch:encoder}\\
S3 embed + blocks & $k3,s2,p1$ & \shape{(B, 64, 160)} $\to$ F3 \shape{(B, 160, 8, 8)} & \ref{ch:encoder}\\
S4 embed + blocks & $k3,s2,p1$ & \shape{(B, 16, 256)} $\to$ F4 \shape{(B, 256, 4, 4)} & \ref{ch:encoder}\\
decoder fuse & $1\times1$ convs, upsample & \shape{(B, 256, 32, 32)} & \ref{ch:decoder}\\
sub-pixel conv & $3\times3$ conv & \shape{(B, 512, 32, 32)} & \ref{ch:decoder}\\
pixel shuffle & factor 4 & \shape{(B, 32, 128, 128)} & \ref{ch:decoder}\\
decompress & \texttt{SpectralDecompressor} & \shape{(B, 165, 128, 128)} & \ref{ch:compressor}\\
denormalise & \texttt{PixelDenormalize} & \shape{(B, 165, 128, 128)} & \ref{ch:bandpipeline}\\
\bottomrule
\end{tabular}
\end{table}

A habit worth adopting: annotate every line of tensor code with its shape, as the repository does
(\file{docs/10-code-style.md}):
\begin{lstlisting}
x = x.flatten(2)          # x: (B, c_out, N)
x = x.transpose(1, 2)     # x: (B, N, c_out)
\end{lstlisting}

%======================================================================
\begin{confusions}
\confusion{Is $N$ the number of pixels?}{Only when you flatten a full-resolution map. After a patch
embedding $N$ counts tokens: at stage~1 each token covers a $4\times4$ block, so $N = (H/4)(W/4)$, one
sixteenth of the pixel count.}
\confusion{Is a reshape the same as a transpose?}{No. A reshape keeps memory order and changes only
the filing system; a transpose changes which number sits where. Converting \shape{(B, C, H, W)} to
\shape{(B, N, C)} needs the transpose (\cref{we:shapes:layout}).}
\confusion{Is a deep-learning convolution the mathematician's convolution?}{It is a cross-correlation:
the kernel is not flipped. Since the kernel is learned, nothing is lost.}
\confusion{Does the stride of a tensor have anything to do with the stride of a convolution?}{No. A
tensor's strides are memory distances between neighbours along each axis; a convolution's stride is
how far the kernel jumps. Same word, two meanings.}
\confusion{Does a bigger image need more convolution parameters?}{No. \Cref{eq:shapes:convcount} has
no $H$ or $W$ in it: the same kernel is reused at every position.}
\confusion{Is a $1\times1$ convolution really a convolution?}{Formally yes, with a window of one pixel.
In practice it is a per-pixel matrix multiply that mixes channels and never looks at neighbours.}
\end{confusions}

\begin{checkyourself}
\item\label{ex:shapes:1} Read \shape{(8, 165, 128, 128)} aloud in plain words. Then convert
\shape{(4, 64, 16, 16)} to the sequence layout: what are $N$ and the feature width?
\item\label{ex:shapes:2} Compute by hand the convolution of
$x = \begin{pmatrix}2&0&1\\1&3&0\\0&1&2\end{pmatrix}$ with
$K = \begin{pmatrix}0&1\\1&-1\end{pmatrix}$, bias $0$, stride~1, no padding.
\item\label{ex:shapes:3} Use \cref{eq:shapes:outsize} to find the output size of a convolution with
$k=7$, $s=4$, $p=3$ on $H=128$ (the stale ``kernel 7'' of \cref{ch:patchembed}), and of $k=3$, $s=2$,
$p=1$ on $H=31$. Is the second equal to $31 /\!/ 2$?
\item\label{ex:shapes:4} How many parameters does \code{nn.Conv2d(64, 160, kernel_size=3)} have? And
\code{nn.Conv2d(32, 165, kernel_size=1)}?
\item\label{ex:shapes:5} A stack of three convolutions has $(k,s) = (3,1)$, $(3,2)$, $(3,1)$. Compute
the receptive field and jump after each layer.
\item\label{ex:shapes:6} Starting from the tensor of \cref{we:shapes:layout} in sequence layout
$\begin{pmatrix}1&5\\2&6\\3&7\\4&8\end{pmatrix}$, apply \code{transpose(1, 2)} then
\code{reshape(1, 2, 2, 2)} and confirm you recover the two channels. What would
\code{reshape(1, 2, 2, 2)} without the transpose give?
\end{checkyourself}

\nextchapter{ch:convae}
